\documentclass[letterpaper,10pt,conference]{ieeeconf}
\IEEEoverridecommandlockouts
\overrideIEEEmargins
\usepackage{amsmath,amssymb,graphicx,booktabs,url}
\title{Motion2Scene: Constructing Training Scenes from Executed Humanoid Motion Contrasts}
\author{Anonymous submission}
\begin{document}
\maketitle
\thispagestyle{empty}
\pagestyle{empty}

\begin{abstract}
Scenes compatible with a motion need not teach a robot when to change its behavior.
We investigate training-scene construction for a humanoid selector choosing between
walking and a fixed crouching command under a frozen controller. Motion2Scene uses
executed command alternatives to propose overhead constraints and paired closed-loop
simulation to label both command outcomes. We compare uniform, target-only, analytic
contrast, and learned contrast construction. Under a nominal-placement contract,
each arm supplies fifteen labeled encounter groups; analytic and learned contrast
construction each acquire nine useful contrasts, while the other arms acquire none.
Their selectors pass 24/36 and 22/36 evaluation conditions, versus 14/36 for either
non-contrast construction. A post-hoc matched-command audit shows that always
requesting the adaptation also passes 24/36. Analytic construction matches this
passage while reducing unnecessary requests from 14/14 to 8/14 both-pass cases;
learned construction makes 6/14 unnecessary requests but misses two useful
adaptations. A retained perturbation-contract failure exposes scarce geometric
eligibility. Full-observation collisions with conflicting labels explain an
irreducible fitting-loss component in the learned corpus. The results support
executed-contrast targeting at this label budget on a finite, partly observed
development bank, without establishing learned-generator superiority, equivalence,
or source-held-out transfer.
\end{abstract}

\section{Introduction}
A humanoid that can walk and crouch must still decide when crouching is useful.
Collecting motions that demonstrate these capabilities does not directly provide
training data for that decision. An overhead beam may permit both commands,
defeat both, or permit only one. Moreover, a beam that clears a crouching reference
may contact the robot while its controller enters the crouch. A scene constructor
must therefore reason about the command the robot executes from its current
history, not only the intended pose.

We ask: \emph{Which constructed scenes teach a frozen humanoid controller's
selector when to adapt?} Our controlled task has one beam family and one decision:
commit to walking or request an existing crouching reference, d040, at 0.30\,s.
The low-level SONIC controller and its command interface remain fixed. Recorded
empty-scene executions provide geometric forecasts; paired obstacle-present
executions determine the labels. A common learner predicts the two physical
outcomes from ideal rays and robot state. Figure~\ref{fig:execution} illustrates
the distinction between the reference and executed command.

This design separates four questions that are easily conflated: whether a
geometric contrast exists, whether it survives execution, whether the selector
can distinguish its observations, and whether the resulting corpus improves
decisions. We vary construction while fixing the downstream learner. The
analytic constructor is a strong comparator within the framework. The learned
proposal is an ablation whose value is measured, rather than an assumed source
of superiority.

Our contributions are (i) a command-bound scene-construction and paired-labeling
protocol retaining all four physical outcomes; (ii) a finite-bank comparison
showing the utility and limits of contrast-targeted training at fifteen groups
per arm, including simple policies; and (iii) diagnostics of perturbation
eligibility and complete-input label aliasing. The perturbation study and the
nominal study remain distinct. The nominal revision followed inspection of
earlier evaluation outcomes; its inference is explicitly conditional on this
development bank. Neither completed simulations nor small descriptive
$p$-values imply deployment robustness or a new-motion-source result.

\begin{figure*}[t]
\centering
\includegraphics[width=0.96\textwidth]{figures/executed-contrast.pdf}
\caption{Why condition construction on execution? One nominal analytic training
encounter on carrier C1, label seed 8722. Left: commanded reference-bank root
height (dashed) and executed height (solid), with the fixed 0.30\,s decision
marked. Right: recorded per-body peak beam force through each command's scoring
horizon. Walking contacts the beam; d040 passes. Root height illustrates tracking
differences and is not itself a clearance certificate. The method uses complete
achieved body trajectories for geometry and paired physical outcomes for labels.}
\label{fig:execution}
\end{figure*}

\section{Related Work}
Learning from Hallucination (LfH) constructs constrained environments from motion
collected in open space~\cite{lfh}. Learning from Learned Hallucination (LfLH)
learns obstacle distributions with an encoder and a planning-based decoder,
then trains reactive planners on rendered observations~\cite{lflh}.
LfH-CP factorizes dynamic-scene generation into critical-point identification
and obstacle trajectories passing through those points, and measures downstream
navigation performance~\cite{lfhcp}. Thus, neither generating environments from
motion nor evaluating their learning utility is new by itself. Our distinction
is the combination of \emph{paired executed humanoid command alternatives},
physical verification under a frozen controller, and a controlled study of
which resulting labels support a fixed adaptation decision. We do not assert
that the paired-label representation is necessary; construction also changes
positive-class coverage and other corpus properties.

HumanoidPF learns whole-body traversal using a humanoid--obstacle potential-field
representation and hybrid scene generation~\cite{humanoidpf}. Perceptive Humanoid
Parkour composes motion-matched skills and trains a perceptive policy for
long-horizon traversal~\cite{php}. These systems address broader capability and
deployment questions. Our idealized sensing and two-command task instead isolate
training-data construction; their system-level performance is not a baseline
for our finite bank. SONIC supplies the inherited tracking capability~\cite{sonic}.
We neither retrain it nor attribute its locomotion skill to Motion2Scene.

Evaluation uncertainty depends on the number and structure of experimental
units~\cite{agarwal}. Matched comparisons do not remove dependence among pairs
sharing a carrier or layout~\cite{eliasziw}. We retain the registered
condition-level analysis, report carrier and layout structure, and distinguish
a failure to detect a difference from equivalence~\cite{lakens}.

\section{Command-Bound Construction and Learning}
\subsection{Task, commands, and physical outcomes}
Let $S(q)$ be a static beam scene, $x_t$ the measured robot state, $h_t$ its
approach history, and $a\in\{0,1\}$ a supported command. The outcome is
$Y_a=Y(S(q),x_t,h_t,a;\xi)$, conditional on controller and physics seed $\xi$.
Command $a=0$ commits to walking. Command $a=1$ requests d040 at 0.30\,s and
returns to walking through the existing 3.3--3.5\,s legality window. This is one
scene-dependent decision. Subsequent return checks use phase and reference
continuity, not a new policy decision. Requests can be denied by the shared
interface and are reported separately from successful adaptations.

The frozen endpoint is contact-qualified passage in the first episode. All
recorded body origins must pass at least 0.10\,m beyond the beam's downstream
face and remain across that boundary and upright for 0.30\,s. Upright means
root height at least 0.50\,m and the vertical component of the projected
upward direction at least 0.50. No per-body beam force may exceed 1\,N before
the qualifying interval finishes. At 50\,Hz, the stability window includes
sixteen samples. Direct 200\,Hz forces are aggregated by the maximum-norm
substep for each body, with synchronization checked against the recorder.

The execution driver requests one 200-control-step reference pass; the records
contain 199 samples over 3.96\,s when no reset occurs. The first reset truncates
the scored episode. Without a qualifying interval, force is evaluated through
that episode's end. Later recovery cannot rescue an earlier failure. Falls,
resets, denied requests, entry/return, and the separate reference-tracking
verdict remain reported diagnostics. The endpoint concerns recorded body
origins and contact samples; it is not native-collider clearance or a
continuous-time collision guarantee.

An \emph{encounter group} contains one scene/carrier/label-seed assignment,
its decision-time observation, and both command outcomes. Its label is
$(Y_0,Y_1)$. Both-pass groups teach walking preference; $(0,1)$ groups provide
useful adaptations; $(1,0)$ groups warn against adapting; and both-fail groups
describe command-bank inadequacy. Missing commands are unknown, not negative
labels. Measurement admission requires matched recorded pre-decision state,
actions, motion tokens, reference history, and complete decision inputs.

\subsection{Geometry is a proposal criterion}
The beam is 0.10\,m deep, 1.20\,m wide, and 0.10\,m thick. Its parameters are
source-route station $s\in[0.10,0.90]$ and underside height
$z\in[1.10,1.45]$\,m. Empty-scene walk and d040 executions supply sampled
achieved body geometry, including the command transition. The reference route
places the beam in world coordinates. Let $c_a(q)$ be the minimum signed
capsule--beam clearance along the recorded command, with positive values
indicating separation. The nominal geometric contrast set is
\begin{equation}
\mathcal C=\{q:c_1(q)\geq m,\ c_0(q)\leq -m\},\qquad m=0.010\;\mathrm m.
\end{equation}
These sampled capsule forecasts are distinct from $Y_a$. Contacts alter
execution, and a geometric margin is not a physical outcome.

The original acceptance contract requires these inequalities at every member
of a fixed 113-offset set in world $x,y,z$, and yaw. Its extrema are
$\pm20$\,mm horizontally, $\pm10$\,mm vertically, and $\pm0.02$\,rad.
It is a \emph{worst-case finite-set} test, not a probabilistic or continuous
uncertainty certificate. The nominal contract evaluates acceptance only at
zero offset. Both retain the original 17-offset search objective; the 17- and
113-offset outcomes are recorded as diagnostics in the nominal study.

For intuition only, suppose a fixed horizontal pose admits one closed interval
of useful heights $[L,U]$. Requiring \emph{every} vertical perturbation in
$[-\epsilon,\epsilon]$ to remain useful gives
\begin{equation}
[L+\epsilon,U-\epsilon],\qquad
w=\max(0,U-L-2\epsilon).
\end{equation}
The interval is empty when the lower endpoint exceeds the upper; equality
leaves a singleton of zero length. This conditional one-dimensional erosion
does not establish interval connectivity, endpoint geometry, or robustness
for our full pose-space test. In particular, our measured 20--27\,mm best
\emph{joint margins} are not interval widths.

\subsection{Four frozen construction arms}
All arms receive the same carrier references and use the same scene domain.
Uniform draws sixteen outputs per carrier without a future-contrast filter.
Analytic construction uses achieved-motion envelopes, a global 136-candidate
search initialized over twenty stations, and distinct-output ranking.
Motion2Scene uses a frozen conditional eight-component proposal model and
bounded pattern search. Its temporal encoder consumes a 64-frame normalized
target-reference sequence and predicts mixture weights and correlated
two-dimensional latent Gaussians. Sigmoid transforms map samples to eight
station intervals and the common height domain. The inherited model is used
with the exact d040-bound input and is not refitted on the evaluation bank.

Each learned output receives seventeen candidate queries: the initial point
and four coordinate directions at four decreasing scales. Search stays within
the original $\pm0.05$ station and $\pm0.03$\,m height box. Target-only uses
the separately frozen no-contrast checkpoint and target-only search and
acceptance; it cannot regain a walking-interference objective through ranking.
This arm is a construction ablation, not a same-corpus single-label control.

All arms reject nonfinite/out-of-domain outputs, within-arm duplicates, reserved
layout neighborhoods, and inadequate pre-decision clearance. The contrast arms
add both margin constraints; target-only adds target clearance. Visibility is
measured but never gates eligibility or selection. Every output and rejected
assignment is retained. Selection precedes physical labeling; no positive-only
refill follows the outcomes.

\subsection{One fixed selector}
The selector receives 214 values: 144 upper/conditional-lower ray features and
70 phase, active-skill, observation-age, gravity, root-velocity, and joint-state
values. Hit distance and hit height are geometric observations; object paths,
source IDs, and true beam parameters are excluded from the feature function.
Robot state can nevertheless encode carrier identity indirectly. Ideal rays
do not constitute a realistic perception-transfer test.

Two logistic heads predict command feasibility. For $G$ complete groups,
\begin{equation}
\mathcal L=\frac{1}{2G}\sum_{g=1}^{G}\sum_{a=0}^{1}
\operatorname{BCE}(p_{ga},Y_{ga})+
0.01\frac{\|W\|_F^2}{428}.
\end{equation}
Input scales are one except the 29 joint velocities, scaled by five. Parameters
start at zero and receive 2,000 full-batch Adam updates at 0.01. Reported BCE
excludes the weight penalty. There are no class weights or test-selected
checkpoints. The linear function is compiled into the existing execution
container and checked for matching probabilities and actions.

At threshold 0.5 the selector prefers walking if its walk head is positive;
otherwise it requests d040 if that head is positive. If neither is positive,
it records a refusal and \emph{continues walking}. Refusal is a classification
output, not a protective stop. Successful adaptation means a requested,
executed d040 command that achieves the passage endpoint; the number of
requests alone does not measure successful adaptation.

\section{Protocols, Provenance, and Metrics}
Three previously inspected development carriers, denoted C1--C3 (source IDs
41001--41003), share the command qualification procedure. Seed 8721 supplies
empty-scene proposal geometry; seed 8722 supplies paired labels. The fixed
evaluation bank crosses stations $\{.35,.45,.55,.65\}$ with underside heights
$\{1.18,1.27,1.36\}$\,m. A \emph{condition} is a carrier, layout, and physics
seed. It is not an independent motion source.

\begin{table*}[t]
\centering\small
\caption{Protocol and deviation summary. Both studies retain proposal seed 8841,
label seed 8722, controller, commands, learner, and scorer.}
\begin{tabular}{p{0.16\textwidth}p{0.36\textwidth}p{0.38\textwidth}}
\toprule
Item & Perturbation study & Nominal study\\
\midrule
Acceptance & 113 fixed offsets; contrast margin 10\,mm & Zero-offset acceptance; same margin; perturbations diagnostic\\
Selection & First six outputs per carrier/arm assigned before audit & First three eligible outputs per carrier/arm selected from sixteen\\
Groups per arm & 18 requested generated + 6 shared backgrounds = 24 & 9 selected generated + the same 6 backgrounds = 15\\
Realized groups & Uniform 24; analytic 7; target-only 22; learned 6 & 15 in every arm; no refill or exclusion after labels\\
Evaluation & 12 layouts $\times$ 2 seeds $\times$ 3 carriers = 72 conditions; six policies & Seed 8511 only: 36 conditions; four new policies; two fixed comparators reused\\
Backgrounds & 108 separate executions: three types, two seeds, three carriers, six policies & Original background labels reused for fitting; parent-model control results reported separately\\
Inference & Unequal acquired data at equal requested slots & Equal group count; outcome-informed revision on overlapping development data\\
\bottomrule
\end{tabular}
\label{tab:protocol}
\end{table*}

Table~\ref{tab:protocol} resolves the 24-versus-15 and 72-versus-36 denominators.
The six shared groups are C1 absent/raised, C2 absent/blocked, and C3
raised/blocked, all at station .60. Each type occurs twice. Shared pairs are
acquired once, even though they enter every arm's fitting set. The nominal
record retains obsolete metadata fields naming six generated slots per
carrier and 72 requested groups; its operative three-eligible quota, code,
and executed assignments agree on 36 generated groups and 72 new commands.
These metadata discrepancies do not justify changing the historical files.

The nominal protocol was recorded after a 366-execution parent snapshot was
available: 61 parent conditions, including 31 of the 36 nominal conditions,
had already been evaluated. The earlier learner revision also used inspected
layout outcomes. Nominal proposals and labels were frozen before their own
execution, but this is not an untouched confirmatory test. The panels reuse
all twelve layouts, three carriers, one evaluation seed, six background labels,
the generator checkpoints, and the controller. Nineteen selected nominal
generated scenes also match previously labeled parent scenes at the same
label seed (nine uniform, nine target-only, one analytic). They were executed
again under the nominal protocol. The two panels are not independent replications.

The registered primary estimand is mean passage within each carrier, averaged
over the three carriers; balanced denominators make it numerically equal to
the pooled fraction. We target \emph{finite-bank performance}, conditional on
the selected training corpora. New-layout population performance and new-source
generalization would require different sampling and holdout designs.

Registered comparisons pair conditions and report every cell of the $2\times2$
table. Exact two-sided binomial tests condition on discordances. Their
independence assumption is not established across reused layouts or carriers;
we retain them as descriptive calculations, not population uncertainty.
Physics repeats stay with their carrier/layout. Deterministic fitting removes
optimizer randomness, not training-corpus variability.

\section{Results and Diagnostics}
\subsection{Acquisition at fifteen groups per arm}
\begin{table*}[t]
\centering\small
\caption{Nominal acquisition and fitted-label composition. Each arm proposes
48 outputs (16 per carrier), selects nine generated groups, and reuses six
background groups. Eligibility is geometric; outcome columns are physical
paired labels over all fifteen groups. Every group has two labels.}
\begin{tabular}{lrrrrrrrr}
\toprule
Construction & Outputs & Eligible & Groups & Both pass & d040 only & Walk only & Neither & BCE \\
\midrule
Uniform & 48 & 48 & 9 + 6 & 7 & 0 & 0 & 8 & 0.001762 \\
Analytic & 48 & 41 & 9 + 6 & 4 & 9 & 0 & 2 & 0.000542 \\
Target-only & 48 & 46 & 9 + 6 & 13 & 0 & 0 & 2 & 0.000175 \\
Motion2Scene & 48 & 33 & 9 + 6 & 4 & 9 & 0 & 2 & 0.139004 \\
\bottomrule
\end{tabular}

\label{tab:acquisition}
\end{table*}
All 72 nominal label commands and 144 nominal policy executions pass measurement
admission. The four primary fits replay with exactly matching weights, biases,
scales, and BCE. Table~\ref{tab:acquisition} separates output count, geometric
eligibility, selected groups, and useful physical contrast. Analytic and learned
construction each provide nine useful contrasts. Uniform and target-only
provide none. Thus, nine useful contrasts out of nine selected generated
groups does not mean nine proposals or equal acquisition computation.

All six nominal predictions hold: each contrast arm supplies at least one
eligible group on each carrier (three selected each); an arm other than analytic
acquires useful contrasts (learned, nine); uniform yields at most one (zero);
target-only is predominantly both-pass (9/9 generated, 13/15 including
backgrounds); precisely the arms with useful labels request d040 at evaluation;
and most nominally eligible outputs fail the 113-offset contrast diagnostic
(167/168). The last count includes untargeted arms; among selected contrast-arm
outputs it is 17/18. No visibility filtering is applied.

\subsection{Passage and the always-adapt alternative}
\begin{table*}[t]
\centering\small
\caption{The same 36 nominal conditions. Requests on both-pass cases are
unnecessary for passage. Both-fail requests and refusals do not avoid failure.
The scripted and privileged geometric comparators reuse their registered
seed-8511 executions. $^{*}$Post-hoc command lookup with eight missing conditions
completed by sixteen separately registered matched executions. The outcome
oracle uses physical labels in hindsight; privileged geometry uses a forecast.}
\begin{tabular}{lrrrrr}
\toprule
Policy & Passage & Requests & Both-pass req. & Both-fail req. & Refusals \\
\midrule
Uniform & 14/36 & 0/36 & 0/14 & 0/12 & 17/36 \\
Analytic & 24/36 & 24/36 & 8/14 & 6/12 & 0/36 \\
Target-only & 14/36 & 0/36 & 0/14 & 0/12 & 0/36 \\
Motion2Scene & 22/36 & 22/36 & 6/14 & 8/12 & 0/36 \\
Scripted rays & 24/36 & 24/36 & 8/14 & 6/12 & 0/36 \\
Privileged geometry & 20/36 & 7/36 & 1/14 & 0/12 & 17/36 \\
Always walk$^{*}$ & 14/36 & 0/36 & 0/14 & 0/12 & 0/36 \\
Always d040$^{*}$ & 24/36 & 36/36 & 14/14 & 12/12 & 0/36 \\
Outcome oracle$^{*}$ & 24/36 & 10/36 & 0/14 & 0/12 & 12/36 \\
\bottomrule
\end{tabular}

\label{tab:comparators}
\end{table*}
Analytic and learned construction pass 24/36 (66.7\%) and 22/36 (61.1\%)
conditions, versus 14/36 (38.9\%) for either non-contrast construction
(Table~\ref{tab:comparators}). Their improvements are 27.8 and 22.2 percentage
points on this bank. Both directions hold on all three carriers. In C1--C3
order, counts are 4/12, 5/12, 5/12 for uniform and target-only; 8/12 on every
carrier for analytic; and 7/12, 7/12, 8/12 for learned construction.
Figure~\ref{fig:paired} preserves layout-level outcomes.

\begin{table}[t]
\centering\footnotesize
\caption{Registered nominal paired comparisons. All four cells are shown.
The condition-level $p$-values assume independent discordances; they are
descriptive because conditions share carriers and layouts. An.: analytic;
Un.: uniform; Tgt.: target-only; M2S: Motion2Scene.}
\setlength{\tabcolsep}{3pt}
\begin{tabular}{lrrrrr}
\toprule
A / B & Both & A only & B only & Neither & $p$ \\
\midrule
M2S / An. & 22 & 0 & 2 & 12 & 0.5 \\
M2S / Un. & 14 & 8 & 0 & 14 & 0.0078125 \\
An. / Un. & 14 & 10 & 0 & 12 & 0.001953125 \\
M2S / Tgt. & 14 & 8 & 0 & 14 & 0.0078125 \\
\bottomrule
\end{tabular}

\label{tab:paired}
\end{table}
No reverse discordances occur against uniform or target-only. Analytic exceeds
learned construction by two conditions (5.6 percentage points), both in the
analytic direction. No statistically significant difference was detected by
the registered test ($p=0.5$); equivalence was not established. A post-hoc
carrier sign calculation has only three positive clusters against uniform
and gives $p=0.25$ if those carriers were independent. We do not use condition
counts to construct narrow new-source confidence intervals.

The strongest simple comparison changes the interpretation. Scripted rays
also pass 24/36, and privileged geometry passes 20/36. Eight nominal conditions
lacked d040 outcomes in the combined records. A separate post-hoc protocol
executed both constant commands on those conditions, retaining the original
settings. All eight repeated walks reproduce passage and recorded pre-decision
history exactly. Across original and audit executions, all 304 available
same-command comparisons match both passage and the recorded state/action/token
trajectory fields. The full 214 inputs, beam, seed, banks, non-policy
configuration, and scoring contract also match across methods. These checks
support command lookup for this one-decision deterministic runtime; they are
not a hidden-state snapshot or an evaluation method for sequential policies.

The completed command table contains fourteen both-pass, ten d040-only,
twelve both-fail, and zero walk-only conditions. Always-d040 therefore passes
24/36, reaching the two-command outcome oracle. Analytic matches this passage
with 24 rather than 36 requests, avoiding six unnecessary adaptations on the
fourteen both-pass cases. Learned construction avoids eight unnecessary
adaptations, but requests d040 on only eight of ten rescue cases and passes
22/36. Its extra selectivity does not imply a superior tradeoff. The scripted
rule makes exactly the same requests and obtains the same outcomes as the
nominal analytic selector on all 36 conditions.

\begin{figure*}[t]
\centering
\includegraphics[width=0.97\textwidth]{figures/paired-outcomes.pdf}
\caption{Where do decisions matter? All nominal carrier/layout conditions,
with the two commands' outcomes completed by the post-hoc matched audit.
Within each carrier, layouts enumerate four increasing stations, each with
three increasing heights. Color reports actual command and passage; a
refusal executes walk and is included in that command's category. The
learned selector's two missed rescues remain visible, as do unnecessary
adaptations and conditions neither command solves.}
\label{fig:paired}
\end{figure*}

Neither contrast selector refuses on the nominal traversal bank. They request
d040 on 6/12 and 8/12 both-fail conditions, respectively. In the parent
background suites, all six policies make zero requests over the 72 absent/raised
executions. All four parent learners and privileged geometry refuse 6/6
blocked cases, whereas scripted rays refuse 0/6; all blocked executions fail.
These are parent-model controls, not evaluations of the new nominal fits.
Successful refusal or safe stopping is not demonstrated. The recorded
nominal analytic and learned selectors each make eighteen successful d040
passages despite making 24 and 22 requests, respectively. Resets occur in
5/36 and 4/36 of their traversal runs; their nominal return counts are 21/24
and 19/22 requests. Passage and complete tracking remain different endpoints.

\subsection{Retaining the failed perturbation contract}
\begin{table}[t]
\centering\small
\caption{The original perturbation study, kept separate. Its equal-24
complete-group goal fails in three arms. These are equal requested slots
with unequal acquired data, not the nominal fifteen-group comparison.}
\begin{tabular}{lrrr}
\toprule
Construction & Groups & Useful & Passage \\
\midrule
Uniform & 24/24 & 0 & 28/72 \\
Analytic & 7/24 & 1 & 39/72 \\
Target-only & 22/24 & 0 & 28/72 \\
Motion2Scene & 6/24 & 0 & 28/72 \\
\bottomrule
\end{tabular}

\label{tab:robust}
\end{table}
The original study completes 82 labeling and 540 evaluation executions:
432 traversal runs over 72 conditions with six policies, plus 108 separate
background controls. Only one generated group, analytic on C1, provides a
useful contrast. The learned arm receives only six shared backgrounds
(Table~\ref{tab:robust}). Its result cannot establish the utility of learned
generated scenes because none entered its corpus.

The prediction that each contrast arm supplies a useful group on every
carrier fails. Uniform's at-most-one prediction holds (0/18 requested), as
does target-only's at-least-twelve both-pass prediction (16/18). All 23
executed generated target-clear predictions achieve passage, satisfying the
zero-false-clear prediction in this selected set; eight receive separate
tracker rejections. The prediction of at least three net improvements for
each contrast arm fails jointly: analytic adds eleven passages, learned adds
zero. All failed predictions remain part of the interpretation.

For the registered paired comparisons, analytic versus uniform has
(both-pass, analytic-only, uniform-only, both-fail) = (28,11,0,33).
Learned versus analytic has (28,0,11,33); learned versus uniform and
target-only each have (28,0,0,44). The eleven-discordance calculation gives
$p=0.0009765625$; zero-discordance comparisons contain no directional
information. Physics repeats remain grouped. Scripted rays pass 46/72,
privileged geometry 39/72. These parent results support the importance of
acquired contrasts, but are not a replication of the nominal comparison.

\begin{figure}[t]
\centering
\includegraphics[width=\columnwidth]{figures/acceptance-envelope.pdf}
\caption{How much eligibility survives the uncertainty requirement? The
existing post-hoc grid diagnostic evaluates 11,421 station/height centers
per carrier. Horizontal half-width is shown; vertical half-width is half
that value and yaw scales to 0.02\,rad. Nominal witnesses 200/168/337 fall
to 2/0/32 at the inherited envelope. These are finite geometric witnesses,
not executed successes or continuous feasibility volumes.}
\label{fig:envelope}
\end{figure}
The stored envelope diagnostic holds recorded execution geometry fixed
(Fig.~\ref{fig:envelope}). Its specified 81-station by 141-height grid has
11,421 centers; the protocol's parenthetical 7,181 total is an arithmetic
error. Nominal contrast support shrinks sharply under the finite offset
test. The observed decrease does not prove continuous infeasibility.
Because the set contains zero offset, excluding nominal failures from
further checks is valid; zero inclusion alone does not prove monotonicity
between differently scaled, non-nested finite sets. The nominal study
addresses a narrower placement-known task and demonstrates no robustness
at the rejected tolerance.

\subsection{Complete-input aliasing and fitting loss}
Among all nominally critical outputs, analytic has decision-time proposal
ray hits in 44/47 scenes, versus 24/34 for learned construction. These
denominators count geometrically critical outputs before common exclusions;
they are not label counts. All nine selected analytic contrasts have ray
hits, compared with five of nine selected learned contrasts. No group is
removed in this post-hoc analysis.

\begin{figure*}[t]
\centering
\includegraphics[width=0.95\textwidth]{figures/full-input-aliasing.pdf}
\caption{Can the complete selector observation distinguish its labels?
Left: two exact full-input equivalence classes in the learned corpus
combine both-pass backgrounds with d040-only contrasts. All 214 values
are identical within each class. Right: fitted mean BCE and the
unrestricted deterministic predictor's empirical BCE infimum from these
collisions. Identical rays alone would not establish this bound. Both
classes share successful d040, so label ambiguity is not an unavoidable
passage error. All original groups are retained.}
\label{fig:alias}
\end{figure*}
The fitting audit checks all 214 features, identical loss reduction,
normalization, regularization, multiplicity, and class balance. Both contrast
arms contain four positive walk labels and thirteen positive d040 labels
over fifteen groups. Analytic BCE is 0.000541579 and learned BCE is 0.139004365,
a ratio of approximately 257. Their separate weight penalties are
0.001360424 and 0.001004991. Exactly reproducing the fixed 2,000-step fits
rules out an accounting mismatch; it does not by itself prove convergence.

In the learned corpus, one complete-input equivalence class contains two
both-pass backgrounds and one d040-only contrast on C1. Another contains
one both-pass background and three d040-only contrasts on C3. Their maximum
absolute full-feature differences are zero. An arbitrary deterministic
predictor must assign a common probability to each identical input. If a
class has $n$ members with positive fraction $r$ in a head, its minimum
summed BCE is $nH(r)$, where $H$ is binary entropy. Summing over classes
and heads, and dividing by thirty labels, gives an empirical infimum of
0.138629 for this corpus. The measured learned BCE is only 0.000375 higher.
Analytic has no conflicting exact-input class and zero observed head
classification errors; learned has two, matching its minimum possible
head-error count. The post-hoc six-, four-, and two-decimal comparisons
remain precision diagnostics, not substitutes for exact collisions.

This is direct evidence for an observation limitation in fitting walk
feasibility, beyond an aggregate correlation with ray visibility. It does
not isolate invisibility as the cause of downstream performance or of
every generator difference. In particular, d040 succeeds for every member
of both conflicting classes: always adapting can resolve their passage
without resolving whether adaptation is necessary. Shared both-pass and
both-fail examples remain useful training information.

\subsection{Dependence on the acquired corpus}
The twenty-four registered parent leave-four-assignment-out refits are
retained, including missing and zero-member folds. Removing the single
analytic contrast produces a model exactly equal to the background-only
learned-arm model. On ninety recorded parent condition inputs, that refit's
requests fall from 29 to zero; the other folds produce 29, 29, 31, 29,
and 29 requests. This recorded-input comparison is a post-hoc diagnostic,
not new physical evaluation of refit policies.

We additionally refit the unchanged nominal learner after removing a
carrier's entire training set, including its backgrounds, and after
removing each contrast group separately. These thirty post-hoc CPU fits
use matched command lookup and never replace the primary selectors.
Leaving out C1/C2/C3 yields 8/12, 6/12, 8/12 passages on the removed
carrier for analytic, and 8/12, 7/12, 7/12 for learned construction.
Uniform and target-only retain 4/12, 5/12, 5/12. Removing one contrast
gives analytic 21--24/36 and learned 22/36 on the full bank. These are
training-set sensitivity measurements, not independent training datasets
or confidence intervals. The generator still includes development-source
information; selector-only removal establishes no generator-level
source-held-out transfer.

\subsection{Resource accounting}
\begin{table}[t]
\centering\small
\caption{Measured rollout-process time, including startup and contention.
The first four rows sum to 6.577196\,h (the historical ``6.6 GPU-hour''
boundary). This is not GPU utilization time or total method training cost.}
\begin{tabular}{llrr}
\toprule
Study & Stage & Runs & Process hours \\
\midrule
Robust & labels & 82 & 0.824108 \\
Robust & evaluation & 540 & 4.310964 \\
Nominal & labels & 72 & 0.480848 \\
Nominal & evaluation & 144 & 0.961276 \\
Post-hoc audit & evaluation & 16 & 0.137223 \\
\bottomrule
\end{tabular}

\label{tab:cost}
\end{table}
Table~\ref{tab:cost} counts every executed label and evaluation in the two
studies, with shared labels and reused comparators charged once. The
sixteen new audit executions add 0.137223\,h separately. The shared
empty-transition bank costs another 0.105115\,h outside that historical
boundary. Controller training, motion/generator pretraining, earlier
experiments, CPU search, and idle scheduling time are excluded. Original
selector-fit CPU wall time was not recorded, so no complete end-to-end
efficiency claim is made.

Nominal construction preparation takes 38.820 CPU wall seconds. Per-arm
recorded construction/search/audit loops take 3.491, 9.420, 8.955, and
15.984\,s for uniform, analytic, target-only, and learned construction.
Their search-clearance query counts are 0, 13,872, 13,872, and 27,744,
respectively, where one query evaluates a candidate, offset, and command.
Common verification and overhead remain additional to that query counter.
Equal-label comparison is therefore not equal-compute comparison. Historical
fresh-source search timing concerns a different reference-only apparatus
and is not evidence of nominal selector efficiency or transfer.

\section{Limitations and Conclusion}
This study uses one controller, one beam family, ideal simulator rays, one
adaptation decision, and three inspected development carriers. Its fixed
bank was partly observed before nominal registration. Generalization to
new layouts drawn from a population, new motion ancestry, realistic
perception, or hardware remains untested. Finite geometry checks and
contact-qualified passage provide neither robust deployment guarantees
nor a protective refusal behavior. Corpus-removal refits expose some
training-data dependence but do not estimate acquisition variability.

At fifteen groups per arm, executed-contrast targeting supplies the
adaptation examples absent from uniform and target-only corpora and
improves measured finite-bank passage. Always-d040 and scripted rays
show the limit of that result: learned scene construction does not
establish superior traversal capability. Analytic reaches the available
command ceiling while avoiding some unnecessary adaptations; learned
construction trades two missed rescues for fewer unnecessary requests.
Preserving the failed perturbation contract and auditing complete-input
aliasing make the framework's restrictions explicit. The contribution
is a measured account of constructing and verifying decision-relevant
training data, not equivalence, learned-generator superiority, or a
claim that paired labels are necessary.

\section*{Acknowledgments}
For the present revision, OpenAI Codex assisted in drafting and revising Sections I--VI, implementing
the submission audit and diagnostic scripts, and producing plotting code
for Figures 1--4 and recorded-state video rendering. The figures plot
recorded or computed scientific data; robot imagery replays measured
states and is not generated by an image model.

\begin{thebibliography}{99}
\bibitem{lfh} X. Xiao, B. Liu, G. Warnell, and P. Stone,
``Toward agile maneuvers in highly constrained spaces: Learning from hallucination,''
\emph{IEEE Robotics and Automation Letters}, vol. 6, no. 2, pp. 1503--1510, 2021.
\bibitem{lflh} Z. Wang, X. Xiao, A. J. Nettekoven, K. Umasankar, A. Singh,
S. Bommakanti, U. Topcu, and P. Stone,
``From agile ground to aerial navigation: Learning from learned hallucination,''
in \emph{IEEE/RSJ International Conference on Intelligent Robots and Systems}, 2021.
\bibitem{lfhcp} S. A. Ghani, K. Lee, and X. Xiao,
``Learning from hallucinating critical points for navigation in dynamic environments,''
arXiv:2509.26513, 2025.
\bibitem{humanoidpf} H. Xue et al.,
``Collision-free humanoid traversal in cluttered indoor scenes,''
arXiv:2601.16035, 2026.
\bibitem{php} Z. Wu et al., ``Perceptive humanoid parkour: Chaining dynamic human
skills via motion matching,'' in \emph{Robotics: Science and Systems}, 2026.
\bibitem{sonic} Z. Luo et al., ``SONIC: Supersizing motion tracking for natural
humanoid whole-body control,'' \emph{Science Robotics}, vol. 11, no. 117,
eaed4592, 2026.
\bibitem{agarwal} R. Agarwal, M. Schwarzer, P. S. Castro, A. C. Courville, and
M. G. Bellemare, ``Deep reinforcement learning at the edge of the statistical
precipice,'' in \emph{Advances in Neural Information Processing Systems}, 2021.
\bibitem{eliasziw} M. Eliasziw and A. Donner, ``Application of the McNemar test to
non-independent matched pair data,'' \emph{Statistics in Medicine}, vol. 10,
no. 12, pp. 1981--1991, 1991.
\bibitem{lakens} D. Lakens, ``Equivalence tests: A practical primer for $t$ tests,
correlations, and meta-analyses,'' \emph{Social Psychological and Personality
Science}, vol. 8, no. 4, pp. 355--362, 2017.
\end{thebibliography}
\end{document}
